\section{第二部分 系统组成及功能说明（五）：固定工位、机械臂与 F407 执行层}

\subsection{固定工位定位}
双 myCobot 固定工位是实验室执行单元，不是普通机械臂展示。它包含 Pi 侧机械臂控制、USB 末端相机、AprilTag/视觉门控、自研夹爪、WorkCockpit、示教位姿和回放证据。完整方案中，两台机械臂分别承担取料、辅助固定、转移和异常接管任务；初赛阶段重点展示 arm01 的视觉门控固定点取袋闭环和单臂冗余接管能力。

\begin{figure}[H]
  \centering
  \includegraphics[width=\linewidth]{fig_xrd_workstation_html.png}
  \caption{固定工位、视觉门控与 F407 末端执行。}
  \label{fig:xrd_workstation}
\end{figure}

\subsection{arm01 视觉门控取袋}
arm01 已完成“OBSERVE 位拍照 -> AI 脑识别 AprilTag -> \texttt{station\_ok=true} -> 回 START -> 固定点取袋/落袋”的闭环。该结果可以写作视觉门控固定点取袋闭环，不能写作泛化视觉伺服抓取。实测记录包括 1280x720 抓帧、AprilTag 36h11 id=8、PICK 误差约 8.5-10.6 mm、START 回位约 1.8-1.9 mm，以及末端 G23 舵机开、闭、搬运保持值。

\subsection{STM32F407 末端执行}
F407 末端执行层已打通串口命令、推杆、舵机、电磁铁、短步进升降诊断和瓶盖铁片吸附/释放动作时序。视频中展示的是安全降级流程：先验证瓶子抓取、保持、释放和下位机动作时序，再在后续加入顶部/底部限位和电磁铁驱动隔离后接入车载 X5 导航到位信号。

\begin{warnbox}
F407 升降轴在负载长行程下仍需补齐限位、驱动保护和隔离设计。因此报告只写“末端执行链路与安全降级演示已跑通”，不写成无人值守完整取放闭环。
\end{warnbox}

\subsection{双臂 v4 蓝图}
v4 方案定义了从配方下发、车载取料、车-工位对接、双臂倒粉、研磨、故障冗余、灌装、送烧结、取烧结样到 XRD/PL 回填的十阶段剧本。报告中将其作为系统蓝图和后续扩展方向；初赛已实测的内容单独列出，避免把蓝图全部写成已完成真机闭环。
